% ID: SEL-D0055-OA-E04
% ESTADO: validada. S056.
\textbf{a)} Sea $p$ la proporción poblacional que vio el programa. Contrastamos $H_0:p\ge0{,}30$ frente a $H_1:p<0{,}30$.

\textbf{b)} La prueba es de cola izquierda y se calcula en la frontera $p_0=0{,}30$, que maximiza la probabilidad de observar proporciones bajas dentro de $H_0$. Hay $150$ éxitos y $350$ fracasos esperados en esa frontera. La aproximación normal da la región crítica $z<z_{0{,}055}\approx-1{,}598$, con
\[
z=\frac{\hat p-0{,}30}{\sqrt{0{,}30\cdot0{,}70/500}}.
\]
En la escala de proporciones, se rechaza por debajo de un umbral de aproximadamente $0{,}26725$.

\textbf{c)} $\hat p=130/500=0{,}26$ y $z\approx-1{,}952$. Está en la región crítica. Rechazamos $H_0$ y concluimos que hay evidencia de que la proporción de espectadores es inferior al $30\%$ al nivel indicado.
